% Lorgat 2020 Conjecture 1: theorem boundary and construction data.

\providecommand{\cO}{\mathcal{O}}
\providecommand{\fg}{\mathfrak{g}}
\providecommand{\Ell}{\operatorname{Ell}}
\providecommand{\wt}{\operatorname{wt}}
\providecommand{\kBKM}{\kappa_{\mathrm{BKM}}}
\providecommand{\kch}{\kappa_{\mathrm{ch}}}
\providecommand{\kcat}{\kappa_{\mathrm{cat}}}
\providecommand{\kfib}{\kappa_{\mathrm{fiber}}}
\providecommand{\gDelta}{\mathfrak{g}_{\Delta_5}}

\section{Lorgat Conjecture 1: theorem boundary and construction data}
\label{sec:lorgat-conjecture-one-theorem-boundary}

Conjecture 1 of Lorgat 2020 relates three mathematical objects:
diagonal-divisor Siegel paramodular forms, reciprocal square roots of
reduced Donaldson-Thomas series of \(K3\times E\) quotients, and
denominator functions of generalized Borcherds-Kac-Moody superalgebras
whose signed root multiplicities are read from twisted K3 elliptic
genera.

The assertion has three layers. The automorphic layer specifies a
Jacobi input, a paramodular group, a character, a divisor, a chamber,
and a Borcherds product. The curve-counting layer identifies the
reduced Donaldson-Thomas scalar in the same normalization. The
Lie-superalgebra layer specifies a root lattice, a positive cone, real
simple roots, a Weyl vector, a parity split for imaginary simple roots,
signed root multiplicities, and a Weyl-Kac-Borcherds denominator
identity. The product side of a Borcherds lift supplies the
automorphic factor. The BKM theorem requires the full root datum.

\subsection{The primitive \texorpdfstring{\(\Delta_5\)}{Delta5} denominator}
\label{subsec:lorgat-conjecture-one-delta-five}

\begin{theorem}[The constructed denominator at \(N=1\)]
\label{thm:lorgat-conjecture-one-delta-five}
Let \(\phi_{0,1}^{\mathrm{EZ}}\) be the Eichler-Zagier weak Jacobi
form of weight \(0\) and index \(1\), normalized by
\[
  \phi_{0,1}^{\mathrm{EZ}}(\tau,z)\big|_{q^0}
  = y^{-1}+10+y .
\]
The K3 elliptic genus is
\[
  \Ell(K3;\tau,z)=2\phi_{0,1}^{\mathrm{EZ}}(\tau,z).
\]
The primitive denominator uses \(\phi_{0,1}^{\mathrm{EZ}}\), hence
\[
  c_1(0)=10,\qquad
  \kBKM(\gDelta)
  =
  \wt(\Delta_5)
  =
  \frac{c_1(0)}{2}
  =
  5 .
\]
Gritsenko, Nikulin, and Borcherds construct the generalized BKM
superalgebra \(\gDelta\) with denominator
\[
  \operatorname{den}(\gDelta)=D_5(2Z)=64^{-1}\Delta_5(2Z),
  \qquad Z\in\mathfrak H_2 .
\]
Thus the BKM-superalgebra layer of Conjecture 1 is a theorem at
\(N=1\).
\end{theorem}

\begin{proof}
The coefficient of \(y^0q^0\) in
\(\phi_{0,1}^{\mathrm{EZ}}\) is \(10\). Borcherds' automorphic-product
theorem gives the product weight as one half of this constant
coefficient. Gritsenko and Nikulin construct the Lorentzian
generalized BKM superalgebra whose denominator is the weight-\(5\)
Siegel paramodular form \(\Delta_5\)
\cite{GritsenkoNikulin1995,GritsenkoNikulin1998,Borcherds1995,Gritsenko1999}.
The local Vol III computation records the monic normalization
\(D_5=64^{-1}\Delta_5\) and the denominator \(D_5(2Z)\). These data
include the real simple roots, Weyl vector, signed multiplicity rule,
and denominator identity in the primitive row.
\end{proof}

The scalar Donaldson-Thomas branch is the square convention:
\[
  \Phi_{10}^{\mathrm{un}}=\Delta_5^2 .
\]
It carries weight \(10\). The primitive BKM denominator carries weight
\(5\). These are distinct normalizations of the same level-one
automorphic object.

\subsection{The five CHL denominator constants}
\label{subsec:lorgat-conjecture-one-chl-constants}

\begin{proposition}[Five-level Borcherds weights]
\label{prop:lorgat-conjecture-one-chl-constants}
For the primitive CHL denominator levels
\[
  N\in\{1,2,3,4,6\},
\]
the Borcherds-input constants and BKM weights are
\[
  (c_1(0),c_2(0),c_3(0),c_4(0),c_6(0))
  =(10,8,6,4,2),
\]
\[
  \bigl(\kBKM(\Phi_1),\kBKM(\Phi_2),\kBKM(\Phi_3),
  \kBKM(\Phi_4),\kBKM(\Phi_6)\bigr)
  =(5,4,3,2,1),
  \qquad
  \kBKM(\Phi_N)=\frac{c_N(0)}{2}.
\]
\end{proposition}

\begin{proof}
This is the CHL ladder recorded in the Vol III universal
Borcherds-weight theorem and in
Remark~\ref{rem:k3e-adjudication-two-conventions}. Gritsenko and
Nikulin give the level-one denominator and the paramodular ladder;
Gritsenko's Borcherds-weight formula identifies the weight with
\(c_N(0)/2\). The Vol III local computation
\texttt{compute/lib/hall\_borcherds\_gate.py} verifies
\(c_1(0)=10\), \(\kBKM(\Delta_5)=5\), and the square weight \(10\).
The table is a weight theorem for automorphic denominators. A BKM
superalgebra theorem at \(N\geq2\) additionally requires the root data
listed in Section~\ref{subsec:lorgat-conjecture-one-higher-level-data}.
\end{proof}

\subsection{\texorpdfstring{\(K3\times E\)}{K3 x E} and the four invariants}
\label{subsec:lorgat-conjecture-one-k3xe-invariants}

For \(Y=K3\times E\), the compact Hodge supertrace is
\[
  \kch(Y)=\sum_{q=0}^{3}(-1)^q h^{0,q}(Y)
  =1-1+1-1=0.
\]
Kunneth multiplicativity gives
\[
  \kcat(Y)=\chi(\cO_{K3})\chi(\cO_E)=2\cdot0=0.
\]
The Heisenberg specialization, the primitive Borcherds denominator, and
the K3 Mukai lattice give the remaining three readings:
\[
  \kappa_{\mathrm{ch}}^{\mathrm{Heis}}(Y)=3,\qquad
  \kBKM(\gDelta)=5,\qquad
  \kfib(Y)=\operatorname{rk}\widetilde H(K3,\mathbb Z)=24.
\]
Thus
\[
  \{\kch,\kappa_{\mathrm{ch}}^{\mathrm{Heis}},\kcat,\kBKM,\kfib\}(Y)
  =
  \{0,3,0,5,24\},
\]
with the four non-redundant construction values
\[
  \{\kcat,\kappa_{\mathrm{ch}}^{\mathrm{Heis}},\kBKM,\kfib\}(Y)
  =
  \{0,3,5,24\}.
\]
The only universal formula in the BKM lane is
\[
  \kBKM(\Phi_N)=c_N(0)/2.
\]
At \(N=1\), this gives \(\kBKM(\Phi_1)=5\), while the compact
chiral-supertrace expression gives
\[
  \kch(K3\times E)+\chi(\cO_E)=0+0=0.
\]
The equality \(5=3+2\) uses the Heisenberg reading \(3\) and the
K3 fibre value \(\chi(\cO_{K3})=2\). The BKM formula remains
\(\kBKM(\Phi_N)=c_N(0)/2\).

\subsection{The Gritsenko and Clery catalogue}
\label{subsec:lorgat-conjecture-one-gc-catalogue}

The Gritsenko and Clery diagonal-divisor catalogue is a separate
eight-row classification of genus-\(2\) paramodular forms with simplest
divisor. Its rows are indexed by triples \((t,N;k)\):
\[
\begin{array}{c|c|c}
\mathrm{row} & (t,N;k) & \wt \\ \hline
1 & (1,1;5) & 5 \\
2 & (2,1;2) & 2 \\
3 & (1,2;3) & 3 \\
4 & (3,1;1) & 1 \\
5 & (1,3;2) & 2 \\
6 & (4,1;\tfrac12) & \tfrac12 \\
7 & (1,4;\tfrac32) & \tfrac32 \\
8 & (2,2;1) & 1
\end{array}
\]
Equivalently,
\[
  \wt=(5,2,3,1,2,\tfrac12,\tfrac32,1),
  \qquad
  c^{\mathrm{GC}}(0)=(10,4,6,2,4,1,3,2),
\]
with
\[
  \wt(F_{t,N})=\frac{c^{\mathrm{GC}}_{t,N}(0)}{2}.
\]
The overlap with the CHL denominator branch is the primitive row
\((t,N;k)=(1,1;5)\), where \(F_{1,1}=\Delta_5\). The remaining rows
live on their own paramodular groups and multiplier systems
\cite{GritsenkoClery2013}. A compact Hall-Borcherds reading of any
catalogue row requires a root lattice, Weyl vector, chamber, reflection
character, parity split, and compact Hall comparison.

\subsection{The ordinary Mathieu frame-shape census}
\label{subsec:lorgat-conjecture-one-mathieu-census}

\begin{proposition}[Ordinary Mathieu frame-shape census]
\label{prop:lorgat-conjecture-one-mathieu-census}
The ordinary \(M_{24}\) twining table is the ATLAS table with \(26\)
conjugacy classes
\[
\begin{gathered}
1A,\ 2A,\ 2B,\ 3A,\ 3B,\ 4A,\ 4B,\ 4C,\ 5A,\ 6A,\ 6B,\ 7A,\ 7B,\\
8A,\ 10A,\ 11A,\ 12A,\ 12B,\ 14A,\ 14B,\ 15A,\ 15B,\ 21A,\ 21B,\ 23A,\ 23B.
\end{gathered}
\]
The frame-shape assignment used for the ordinary twining layer includes
\[
\begin{array}{c|c|c}
\mathrm{class} & \mathrm{order} & \mathrm{frame\ shape} \\ \hline
3B & 3 & 3^8 \\
4A & 4 & 2^4\,4^4 \\
4B & 4 & 1^4\,2^2\,4^4 \\
6B & 6 & 6^4 \\
10A & 10 & 2^2\,10^2 \\
12A & 12 & 2\,4\,6\,12 \\
12B & 12 & 12^2
\end{array}
\]
Every row satisfies \(\sum_i i a_i=24\) for frame shape
\(\prod_i i^{a_i}\). The \(26\)-class table is the ordinary Mathieu
table; quotienting by power conjugacy gives a different census.
\end{proposition}

\begin{proof}
The ATLAS \(24\)-point permutation character gives the cycle shapes of
the \(M_{24}\) classes \cite{ConwayCurtisNortonParkerWilson1985}.
Gaberdiel, Hohenegger, and Volpato give the corresponding Mathieu
twining data \cite{GaberdielHoheneggerVolpato2010}. The local executable
table \texttt{compute/lib/mathieu\_moonshine\_yangian.py} stores all
\(26\) rows and asserts \(\sum_i i a_i=24\) and trace consistency at
module import. The targeted tests in
\texttt{compute/tests/test\_mathieu\_moonshine\_yangian.py} verify the
\(26\)-class count, frame-shape totals, traces, and the displayed
special rows.
\end{proof}

\subsection{Higher-level BKM construction data}
\label{subsec:lorgat-conjecture-one-higher-level-data}

For every row \(N\geq2\) in Conjecture 1, a BKM-superalgebra theorem
requires explicit data
\[
\begin{array}{ll}
\text{root lattice} & \Lambda_N\ \text{with positive cone and real simple roots},\\
\text{Weyl vector} & \rho_N\ \text{matching the product prefactor},\\
\text{parity split} &
  \Delta^{\mathrm{im}}_{N,\bar0}\sqcup\Delta^{\mathrm{im}}_{N,\bar1},\\
\text{multiplicities} &
  \operatorname{sdim}\,\fg_{N,\alpha}
  =c_N\bigl((\alpha,\alpha)/2,\ell(\alpha)\bigr),
\end{array}
\]
together with the Weyl-Kac-Borcherds identity
\[
  \Phi_N(Z)=e^{-(\rho_N,Z)}
  \prod_{\alpha\in\Delta_N^+}
  (1-e^{-(\alpha,Z)})^{\operatorname{sdim}\,\fg_{N,\alpha}}
\]
\[
  =
  \sum_{w\in W_N}\det(w)\,
  w\!\left(e^{-\rho_N}
  \sum_{a\in\mathcal I_N}\epsilon_N(a)e^{-a}\right).
\]
The weight identity \(\kBKM(\Phi_N)=c_N(0)/2\) is a necessary
automorphic invariant. The theorem-level BKM assertion is the
construction of the displayed root datum and denominator identity.

\subsection{The theorem boundary}
\label{subsec:lorgat-conjecture-one-theorem-boundary-result}

The \(N=1\) row is constructed: \(\Delta_5\), the weight
\[
  \kBKM(\gDelta)=5,
\]
the denominator \(64^{-1}\Delta_5(2Z)\), and the generalized BKM
superalgebra \(\gDelta\). The reduced \(K3\times E\) scalar lies on the
square branch \(\Delta_5^2\).

For \(N\in\{2,3,4,6\}\), the constants
\[
  c_N(0)=(8,6,4,2),
  \qquad
  \kBKM(\Phi_N)=(4,3,2,1)
\]
are automorphic denominator weights in the CHL normalization. The
remaining mathematical work is the construction of the corresponding
BKM root lattices, Weyl vectors, parity rules, signed multiplicities,
and Weyl-sum identities, plus the scalar comparison with the reduced
Donaldson-Thomas series in the same normalization.

\paragraph{Primary anchors.}
Lorgat 2020;
Gritsenko and Nikulin~\cite{GritsenkoNikulin1995,GritsenkoNikulin1998};
Borcherds~\cite{Borcherds1995,Borcherds1998};
Gritsenko~\cite{Gritsenko1999};
Gritsenko and Clery~\cite{GritsenkoClery2013};
ATLAS~\cite{ConwayCurtisNortonParkerWilson1985};
Gaberdiel, Hohenegger, and Volpato~\cite{GaberdielHoheneggerVolpato2010}.
